\subsection{Infinite non-prosoluble profinite CN-groups: Problem 20.90}
\label{cand:20.90}

The question asks for an infinite, topologically finitely generated
profinite group with pronilpotent centralizers of nonidentity elements,
which is not prosoluble. The archive gives two constructions, counted
as one answer. The first has abelian centralizers.

Let \(R=\F_4[[t]]\), with \(\omega^2+\omega+1=0\), and put
\[
 X=\begin{pmatrix}1&1\\0&1\end{pmatrix},\qquad
 Y=\begin{pmatrix}0&1\\1&\omega+t\end{pmatrix},\qquad
 G=\overline{\langle X,Y\rangle}\leq\SL_2(R).
\]
The congruence topology makes \(\SL_2(R)\) profinite, hence so
is its closed subgroup \(G\), which is two-generated by definition.
Writing \(z=\omega+t\), define
\(P_0=0,P_1=1,P_{n+1}=zP_n+P_{n-1}\). Then
\[
 Y^n=\begin{pmatrix}P_{n-1}&P_n\\P_n&P_{n+1}\end{pmatrix}.
\]
The polynomial \(P_{n+1}\) is monic of degree \(n\) for \(n\geq1\).
Thus no positive power of \(Y\) is the identity and \(G\) is infinite.

Reduction modulo \(t\) is onto \(\SL_2(\F_4)\). To check this
directly, write \(x,y\) for the reduced matrices and
\(U(a)=\left(\begin{smallmatrix}1&a\\0&1\end{smallmatrix}\right)\),
\(L(a)=\left(\begin{smallmatrix}1&0\\a&1\end{smallmatrix}\right)\).
Matrix multiplication gives
\[
 yxy^{-1}=L(1),\qquad
 d=y^{-1}xy^{-2}=\operatorname{diag}(\omega,\omega^2).
\]
Conjugation by \(d\) multiplies the parameters of \(U\) and \(L\)
by \(\omega^2\) and \(\omega\), respectively. Addition of
parameters therefore supplies all elementary matrices, which
generate \(\SL_2(\F_4)\) by row reduction. The commutators
\[
 dU(a)d^{-1}U(a)^{-1}=U(\omega a),\qquad
 dL(a)d^{-1}L(a)^{-1}=L(\omega^2a)
\]
show that this nontrivial finite group is perfect. Hence the
continuous quotient is nonsoluble, and \(G\) is not prosoluble.
This quotient is the usual group \(A_5\).

Regard \(G\) inside \(\SL_2(\F_4((t)))\). In characteristic
two, a scalar determinant-one matrix is the identity.
Every nonidentity \(g\in G\) is therefore nonscalar.
The cyclic-vector argument of Section~\ref{cand:10.35}
identifies its matrix centralizer with the commutative algebra
\(\F_4((t))[g]\). Thus \(C_G(g)\) is abelian.
It is closed, hence an abelian profinite group and in
particular pronilpotent. This field-level centralizer
argument is classical; its characteristic-two form appears
in \cite{fine-ct}. We do not claim it as new.

Thus the example is a CA-group (its nonidentity centralizers are
abelian), hence a CN-group (those centralizers are pronilpotent).
It is not necessary that finite quotients be CN-groups.
Indeed, the archive's image modulo \(t^2\) has order \(240\)
and contains the nonidentity scalar \((1+t)I\), whose
centralizer is the whole nonsoluble image. The theorem
concerns centralizers in the inverse-limit group itself.

\paragraph{A characteristic-zero companion.}
Let \(O=\Z_2[\omega]\subset K=\Q_2(\omega)\) with the same
quadratic relation, and take
\(H=\SL_2(O)/\{I,-I\}\).
The four matrices \(U(1),U(\omega),L(1),L(\omega)\)
topologically generate \(\SL_2(O)\): their powers generate
parameters in the dense subring \(\Z+\Z\omega\), and
elementary row reduction works over the local ring \(O\).
The image of \(U(1)\) has infinite order, and reduction
gives the nonsoluble quotient \(\SL_2(\F_4)\).

For a nonidentity \(h\in H\), a lift \(g\in\SL_2(O)\)
is nonscalar. Its strict centralizer \(C_{\SL_2(O)}(g)\)
is abelian by the same matrix lemma. The preimage of
\(C_H(h)\) consists of \(s\) with \(sgs^{-1}=g\) or \(-g\);
the possible sign is a homomorphism with that strict
centralizer as kernel. If \(\tr g\ne0\), only the
positive sign occurs, so the projective centralizer is
abelian.

If \(\tr g=0\), its reduction \(\bar g\) cannot be the
identity. Otherwise \(g=I+2A\) gives \(\tr A=-1\), and
\[
 1=\det(I+2A)=1+2\tr A+4\det A
\]
would give \(2\det A=1\) in \(O\).
Thus \(\bar g\) is a nonidentity unipotent of
\(\SL_2(\F_4)\), with centralizer of order four
(the matrices \(U(a)\) after a change of basis).
The strict centralizer has reduction a \(2\)-group and
kernel in the pro-\(2\) congruence subgroup, whose
successive quotients embed in \(M_2(\F_4)^+\).
It is pro-\(2\). The extension by a sign group and
subsequent quotient by \(\{I,-I\}\) are also pro-\(2\).
Every projective centralizer is therefore pronilpotent.
Sources: \artifact{research/20.90-proof.md} and
\artifact{research/20.90-2adic-proof.md}.

